\documentclass[11pt,letterpaper]{article}
\usepackage[T1]{fontenc}
\usepackage[utf8]{inputenc}
\usepackage{lmodern,microtype}
\usepackage[margin=0.79in,headheight=15pt]{geometry}
\usepackage{parskip,enumitem,longtable,tabularx,booktabs,array,xcolor,fancyhdr,hyperref}
\definecolor{navy}{HTML}{17314C}
\definecolor{teal}{HTML}{087F78}
\hypersetup{colorlinks=true,linkcolor=navy,urlcolor=teal}
\setlength{\emergencystretch}{3em}
\setlist[itemize]{leftmargin=*,nosep,topsep=5pt}
\setlist[enumerate]{leftmargin=*,itemsep=3pt}
\renewcommand{\arraystretch}{1.22}
\pagestyle{fancy}\fancyhf{}
\fancyhead[L]{\small Enterprise Data Modeling}\fancyhead[R]{\small Documentation Suite}
\fancyfoot[C]{\thepage}
\newcommand{\srcbook}{\textit{The Data Model Resource Book}, revised edition, Volume 2}
\newcommand{\ResponseLine}{\par\vspace{1.05em}\noindent\rule{\linewidth}{0.4pt}\par}
\newcommand{\ResponseArea}{\ResponseLine\ResponseLine\ResponseLine}
\begin{document}
\begin{center}
{\LARGE\bfseries\color{navy} Chapter 5: Insurance}\par
\vspace{0.25em}{\large Study guide | concept, critique, design, verify}\par
\vspace{0.3em}{\small Original engineering study material \textbullet{} 2026 edition}
\end{center}
\vspace{0.8em}
\section{Learning outcomes}
\begin{itemize}
\item Explain how the Chapter 5 domain uses reusable enterprise patterns without assuming that the reference model is a deployable schema.
\item Define at least five entities or associations relevant to insurance and state identity, lifecycle, and cardinality assumptions.
\item Create a conceptual view, a logical ERD, and a small physical schema with explicit business rules.
\item Defend a domain-specific modeling decision using a concrete positive and negative example.
\item Trace each model element to source material, domain requirements, and verification evidence.
\end{itemize}

\section{Chapter orientation}
Primary reading: Chapter 5, Insurance, printed pages 171--246, in the supplied revised Volume 2. The chapter investigates insurance products, coverages, rates, policy periods, insured parties, agreements, claims, loss events, loss reserves, adjudication and payments. The book provides industry modeling context; the implementation tasks and review criteria in this guide are independently authored supplementary practice.

Universal models are reusable hypotheses, not mandatory canonical physical schemas. Distinguish a named business entity from its role in a transaction, and distinguish intended states from observed events. Preserve the exact vocabulary, effective time, optionality and exceptions discovered with domain experts.


\section{Patterns to recognize}
\begin{tabularx}{\linewidth}{@{}p{0.22\linewidth} p{0.37\linewidth} p{0.37\linewidth}@{}}
\toprule
\textbf{Modeling pattern} & \textbf{Design obligation} & \textbf{Illustrative example} \\ 
\midrule
Coverage & Separate covered peril, limit and deductible & Automobile liability with policy-specific limits \\ 
Policy term & Retain versions and effective/expiration dates & Midterm endorsement changes coverage \\ 
Insured role & Associate people, organizations and insured assets & Named insured and covered driver \\ 
Claim / loss & Keep reported event and adjudication separate & Accident generates multiple claim lines \\ 
Financial settlement & Link reserves, payments and recoveries & Partial settlement followed by subrogation \\ 
\bottomrule
\end{tabularx}

\section{Guided domain scenario}
An insurer manages endorsements, overlapping insured parties, claim reserves and payments across policy periods.

\begin{itemize}
\item List actors, business events and system boundaries before naming tables.
\item Propose a stable identity scheme for each independently managed entity.
\item Model mandatory and optional relationships, including allowed multiplicities and time bounds.
\item Show one rule that can be enforced by a database constraint and one that requires a transactional/application policy.
\end{itemize}

\section{Incorrect designs to challenge}
\begin{itemize}
\item Overwriting past policy coverage
\item Assuming a loss equals an approved claim
\item Omitting benefit limits and currency
\item Combining reserves with actual cash movements
\end{itemize}
For each anti-pattern, document the failure mode, an example input that exposes it, and one structural or operational mitigation. Defend why a different candidate model would be less reliable for the stated requirements.


\section{Laboratory tasks}
\begin{enumerate}
\item Construct a context glossary with at least twelve business terms and explicitly define six entity types.
\item Draw a conceptual ERD with named relationship verbs and optionality; then create a physical logical-to-SQL mapping.
\item Model one temporal relationship with effective\_from and effective\_to semantics; specify whether intervals can overlap.
\item Write at least three valid scenarios and three invalid scenarios; map each to the appropriate test or constraint.
\item Present a ten-minute model review focused on identity, referential integrity, temporal validity and ownership.
\end{enumerate}

\section{Knowledge check}
\begin{itemize}
\item How is coverage eligibility evaluated at loss time?
\item Can claims span multiple policy coverages?
\item What is the audit trail for adjustment decisions?
\end{itemize}
Answer each question with a concrete record-level example and identify which facts are explicit in the reference versus assumptions introduced for the exercise.


\section{Completion evidence}
Submit the glossary, one-page model rationale, annotated ERD, physical schema, sample data, test transcript and a completed schema assessment scorecard. A passing review requires consistency among the written rules, diagram multiplicities and executable constraints.


\vfill
\hrule
{\footnotesize\textbf{Provenance.} Len Silverston, The Data Model Resource Book, Revised Edition, Volume 2: A Library of Universal Data Models by Industry Types (Wiley, 2001). Chapter 5, printed pp. 171--246; exercises and modern implementation practices are original supplements.}
\end{document}
